\begin{theorem}[Uniqueness of the pair in the Gaussian obstacle decomposition]
Let $F \in L^2(\gamma_n)$ be nonnegative $\gamma_n$-a.e., put
$m = \int F \, d\gamma_n$, and let $a > m$.  Write $B = L^{1/2}$ for the
operator of \noderef{FracPow}, whose domain $\operatorname{Dom}(B)$ is the case
$s = 1/2$ of \noderef{FracPowDomain}.  Suppose that for each $i \in \{1,2\}$ the
pair $(u_i, g_i)$ with $u_i \in \operatorname{Dom}(B)$ and $g_i \in L^2(\gamma_n)$
satisfies
\[
  u_i \geq 0 \ \text{a.e.},
  \qquad
  F = g_i + B u_i \ \text{ in } L^2(\gamma_n),
  \qquad
  g_i \leq a \ \text{a.e.},
  \qquad
  g_i = a \ \text{a.e. on } \{u_i > 0\} .
\]
Then $u_1 = u_2$ and $g_1 = g_2$ in $L^2(\gamma_n)$.

The hypothesis list is exactly the reduced one from which the reference paper
derives the uniqueness assertion of its obstacle-decomposition theorem: the
mass identity $\int g_i \, d\gamma_n = m$, the lower bound
$\min\{F,a\} \leq g_i$, the measure bound $a\,\gamma_n(\{u_i>0\}) \leq m$ and the
gradient-support identity $Du_i = 0$ a.e.\ off $\{u_i > 0\}$ are deliberately
\emph{not} assumed here, so the conclusion applies to every pair meeting the
four displayed conditions, in particular to any two pairs satisfying the full
conclusion list of that theorem.  The ambient hypotheses $F \geq 0$ a.e.,
$m = \int F$ and $a > m$ are those of the theorem; the first is carried along as
part of the ambient setting and the last two are what make the strict height
inequality $a - m > 0$ available.

Finally, a remark on the way the reference paper phrases the potential's
regularity.  Its decomposition theorem produces $u$ in
``$\operatorname{Dom}(B) = W^{1,2}(\gamma_n)$''.  That equality is the background
identity recalled in the paper's section on the Ornstein--Uhlenbeck operator,
obtained there from the square-root part of the representation theorem; it is
not an assertion that the decomposition theorem proves, and the decomposition
argument consumes only the forward inclusion
$\operatorname{Dom}(L^{1/2}) \subseteq W^{1,2}(\gamma_n)$ together with the norm
identity $\|L^{1/2}v\|_2 = \|Dv\|_2$.  Accordingly the present statement places
the potentials in $\operatorname{Dom}(B) = \operatorname{Dom}(L^{1/2})$ only.
\end{theorem}
\begin{proof}
Throughout, $\mathcal{E}(\cdot)$ is the energy of \noderef{Energy},
$\mathcal{E}(\cdot,\cdot)$ the form of \noderef{EnergyBil}, $\mathcal{K}$ the
cone of \noderef{ObstacleCone}, and $J$ the functional of
\noderef{ObstacleFunctional} built from the given $F$ and $a$, so that
$J(z) = \tfrac12\mathcal{E}(z) + \int (a - F)\,z \, d\gamma_n$.  By
\noderef{GaussIsProb} the measure $\gamma_n$ is a probability measure, so the
constant function $a\mathbf{1}$ lies in $L^2(\gamma_n)$ and
$a\mathbf{1} - F \in L^2(\gamma_n)$.

Fix $i \in \{1,2\}$ and abbreviate $u = u_i$, $g = g_i$.  Put
\[
  \eta = a\mathbf{1} - g \in L^2(\gamma_n) .
\]
Steps 1--6 show that $u$ minimizes $J$ over $\mathcal{K}$; since $i$ is
arbitrary, this holds for both $u_1$ and $u_2$.

\emph{Step 1: $u \in \mathcal{K}$.}
By hypothesis $u \in \operatorname{Dom}(L^{1/2})$, so
\noderef{EnergyFiniteOfSqrtDomain} gives $\mathcal{E}(u) < \infty$.  Together
with $u \geq 0$ $\gamma_n$-a.e., the defining description of the cone in
\noderef{ObstacleCone} gives $u \in \mathcal{K}$.

\emph{Step 2: $\eta \geq 0$ a.e.\ and $\eta\,u = 0$ a.e.}
Since $g \leq a$ a.e., we get $\eta = a - g \geq 0$ a.e.  For the product, let
$N_1, N_2, N_3$ be $\gamma_n$-null sets off which, respectively, $u \geq 0$,
$g \leq a$, and the implication $u(x) > 0 \Rightarrow g(x) = a$ holds.  Their
union $N$ is null by subadditivity of a measure (shared context).  Let
$x \notin N$.  If $u(x) = 0$ then $\eta(x)\,u(x) = 0$.  Otherwise $u(x) \neq 0$,
and $u(x) \geq 0$ gives $u(x) > 0$; then $g(x) = a$, so $\eta(x) = 0$ and again
$\eta(x)\,u(x) = 0$.  Hence $\eta\,u = 0$ $\gamma_n$-a.e.

\emph{Step 3: $Bu = F - g$ and the multiplier identity.}
The hypothesis $F = g + Bu$ is an identity in $L^2(\gamma_n)$, so
$Bu = F - g$ there.  Adding $a\mathbf{1} - F$ to both sides,
\begin{equation*}
  Bu + a\mathbf{1} - F = a\mathbf{1} - g = \eta \quad\text{in } L^2(\gamma_n).
  \tag{$\dagger$}
\end{equation*}

\emph{Step 4: the quadratic expansion of $J$ at $u$.}
Let $v \in \mathcal{K}$ be arbitrary and put $w = v - u \in L^2(\gamma_n)$.  By
Step 1 and \noderef{ObstacleCone} we have $\mathcal{E}(u) < \infty$ and
$\mathcal{E}(v) < \infty$.  Clause (4) of \noderef{EnergyFormBilinear}, applied
first with $v_1 = u$ and $c = -1$, gives $\mathcal{E}((-1)\,u) < \infty$; applied
again to the two finite-energy elements $v$ and $(-1)u$ it gives
$\mathcal{E}(w) = \mathcal{E}(v + (-1)u) < \infty$.

All three of $\mathcal{E}(u), \mathcal{E}(w), \mathcal{E}(v)$ being finite, with
$v = u + w$, clause (3) of \noderef{EnergyFormBilinear} applies to every pair
drawn from $\{u, w, v\}$.  Additivity in the first argument, with second
argument $v = u + w$, gives
$\mathcal{E}(v, v) = \mathcal{E}(u, v) + \mathcal{E}(w, v)$.  Symmetry and
additivity again give
$\mathcal{E}(u, v) = \mathcal{E}(v, u) = \mathcal{E}(u,u) + \mathcal{E}(w,u)$ and
$\mathcal{E}(w, v) = \mathcal{E}(v, w) = \mathcal{E}(u,w) + \mathcal{E}(w,w)$.
Using symmetry once more, $\mathcal{E}(w,u) = \mathcal{E}(u,w)$, so
\[
  \mathcal{E}(v,v) = \mathcal{E}(u,u) + 2\,\mathcal{E}(u,w) + \mathcal{E}(w,w).
\]
By clause (2) of \noderef{EnergyFormBilinear}, $\mathcal{E}(z,z) = \mathcal{E}(z)$
as real numbers for each of $z \in \{u, w, v\}$, all three having finite energy.
Hence
\begin{equation*}
  \tfrac12\bigl(\mathcal{E}(v) - \mathcal{E}(u)\bigr)
    = \mathcal{E}(u, w) + \tfrac12\,\mathcal{E}(w). \tag{3}
\end{equation*}

The linear part of $J$ is handled directly: $a\mathbf{1} - F$, $u$ and $v$ all
lie in $L^2(\gamma_n)$, so the products $(a-F)u$ and $(a-F)v$ are
$\gamma_n$-integrable by the integral form of Cauchy--Schwarz (shared context),
and by linearity of the integral
$\int (a-F)v - \int(a-F)u = \int (a-F)\,w$.  Combining this with $(3)$ and the
formula for $J$ in \noderef{ObstacleFunctional},
\begin{equation*}
  J(v) - J(u)
    = \tfrac12\,\mathcal{E}(w) + \mathcal{E}(u,w) + \int (a-F)\,w \, d\gamma_n .
  \tag{4}
\end{equation*}

\emph{Step 5: the variational inequality.}
Since $u \in \operatorname{Dom}(L^{1/2})$ and $w \in L^2(\gamma_n)$ is arbitrary,
\noderef{EnergyBilIsOperator} applies and gives
$\mathcal{E}(u, w) = \langle Bu, w\rangle$, and
$\langle Bu, w \rangle = \int (Bu)\,w \, d\gamma_n$ by
\noderef{L2InnerGaussIntegral}.  The products $(Bu)w$ and $(a-F)w$ are
$\gamma_n$-integrable by Cauchy--Schwarz (shared context), all factors lying in
$L^2(\gamma_n)$, so by linearity of the integral and $(\dagger)$,
\[
  \mathcal{E}(u,w) + \int (a-F)\,w \, d\gamma_n
    = \int \bigl(Bu + a - F\bigr)\,w \, d\gamma_n
    = \int \eta\, w \, d\gamma_n .
\]
Both $\eta v$ and $\eta u$ are $\gamma_n$-integrable by Cauchy--Schwarz, so
$\int \eta\,w = \int \eta\,v - \int \eta\,u$, and $\int \eta\,u = 0$ because
$\eta u = 0$ a.e.\ by Step 2.  Finally $\eta \geq 0$ a.e.\ by Step 2 and
$v \geq 0$ a.e.\ because $v \in \mathcal{K}$ (\noderef{ObstacleCone}), so
$\eta v \geq 0$ a.e.\ and therefore $\int \eta v \, d\gamma_n \geq 0$.  Altogether
\begin{equation*}
  \mathcal{E}(u,w) + \int (a-F)\,w \, d\gamma_n
    = \int \eta\, v \, d\gamma_n \geq 0 . \tag{5}
\end{equation*}

\emph{Step 6: $u$ minimizes $J$ over $\mathcal{K}$.}
By \noderef{Energy} the energy is a sum of nonnegative terms, so
$\mathcal{E}(w) \geq 0$, and it is a real number by Step 4.  Substituting $(5)$
into $(4)$,
\[
  J(v) - J(u) = \tfrac12\,\mathcal{E}(w) + \int \eta\,v \, d\gamma_n \geq 0
\]
for every $v \in \mathcal{K}$.  That is, $J(u) \leq J(v)$ for all
$v \in \mathcal{K}$.

\emph{Step 7: conclusion.}
Steps 1 and 6 hold for $i = 1$ and $i = 2$, so $u_1, u_2 \in \mathcal{K}$ and
each minimizes $J$ over $\mathcal{K}$.  The hypotheses of
\noderef{ObstacleMinimizerUnique} are therefore met with this $F$, this
$m = \int F \, d\gamma_n$ and this $a > m$, and that lemma gives $u_1 = u_2$.
By Step 3 applied to each index, $g_i = F - B u_i$, whence
\[
  g_1 = F - Bu_1 = F - Bu_2 = g_2 .
\]
This is the asserted uniqueness of the pair.
\end{proof}
